% F04：本书原创矢量示意；依据所在节的定义、证明或明确模型。
\begin{figure}[htbp]
\centering
\begin{tikzpicture}[x=1cm,y=1cm,>=stealth,every node/.style={font=\small},line width=.65pt]

\foreach \s in {0,7}{\draw[->,gray] (\s,0)--(\s+4.5,0) node[right]{\(x\)};
\draw[->,gray] (\s,0)--(\s,2.8) node[above]{\(y\)};
\fill[AnalysisBlue!12] (\s,0)--(\s+2,2)--(\s+4,0)--cycle;
\draw[AnalysisBlue,thick] (\s,0)--(\s+2,2)--(\s+4,0);
\node[below] at (\s+4,0){\(1\)};\node[left] at (\s,2){\(1/2\)};}
\draw[orange!80!black,very thick] (1,0)--(1,1);
\draw[orange!80!black,very thick] (3,0)--(3,1);
\node at (2,3.45){竖直截面：\(\min(x,1-x)\)};
\draw[orange!80!black,very thick] (8,1)--(10,1);
\node at (9,3.45){水平截面：\(1-2y\)};
\node at (5.7,-.8){\(E=\{(x,y):0<x<1,\ 0<y<\min(x,1-x)\}\)};

\end{tikzpicture}
\caption[同一积分域的两种切片]{同一积分域的两种切片。同一三角域按两种方向切片，截面长度的积分均为四分之一。边界为零测，不影响本题面积。}
\label{b1:fig:visual-f04}
\end{figure}
